\section{Realizing the decorated lattices by actual unitary adeles}
\label{sec:decorated-realization}

\status{DR1 proves the needed local Hermitian normal form, including
residue characteristic two, from the literal unramified CM-field
condition. DR2 constructs a unitary finite-adelic matrix carrying the
standard lattice to any supplied self-dual lattice. DR3 installs the
prescribed residue functionals, auxiliary reduction bases and rational
frames, and proves every field of the source's \texttt{IsPtOf}.
Thus the initial decorated-point realizability used in FC and NB is
now supplied under the original CM-field and valid-decoration hypotheses.
The existence of those arithmetic host and column-law data remains a
separate input.}

\paragraph{Conventions and exact arithmetic input.}
Let $E/F$ be the supplied quadratic CM extension, with involution
$a\mapsto\bar a$ and fixed field $F$. Put $\mathcal O=\mathcal O_E$
and use $h(x,y)=\bar x^{\mathsf T}y$ on $E^n$.
The literal \texttt{CMData.cmUnram} condition includes
\[
 \text{for every maximal }P\subset\mathcal O\text{ there is }
 \theta\in\mathcal O\text{ with }\theta-\bar\theta\notin P.
 \tag{DR.1}
\]
This is the local arithmetic assertion used below. It is part of the
supplied CM data; we do not infer it from the name of that record.
Its other clauses identify the involution with complex conjugation
under the embeddings. We use \texttt{CMData.fixed} and
$[E:F]=2$ to identify its fixed field exactly with the given $F$.

For a finite prime $v$ of $F$, write
\[
 K=F_v,\quad R=\mathcal O_{F_v},\quad
 B=E\otimes_F K,\quad
 A=\prod_{P\mid v}\mathcal O_{E_P}\subset B.
 \tag{DR.2}
\]
Thus $A$ is the integral algebra in the completed CM algebra, with
the same involution, not an arbitrarily chosen quadratic extension.
The completion decomposes over the primes above $v$ as in the local
identification construction. Its fixed algebra is $K$, and its fixed
integral subring is $R$. For the first assertion, take any nonfixed
$\theta\in E$: the basis $1,\theta$ remains a $K$-basis after
extension, and comparing the two coefficients of a fixed element
gives the assertion. The second follows by intersecting with the
integral ring, since a scalar in $K$ is integral in all these finite
extensions exactly when its $v$-valuation is nonnegative.

\paragraph{DR1: local orthonormal bases, including the dyadic case.}
First derive the precise local algebra from (DR.1). Choose a prime
$P\mid v$ and its $\theta$ from that condition. If there is another
prime above $v$, it is $\bar P$, since $E/F$ is quadratic Galois.
The element $\delta=\theta-\bar\theta$ is a unit there as well:
$\bar\delta=-\delta$, so membership at $\bar P$ would imply
membership at $P$. Consequently $\delta\in A^\times$.
For every $a\in A$, the elements
\[
 b=\frac{a-\bar a}{\theta-\bar\theta},\qquad
 c=a-b\theta                                             \tag{DR.3}
\]
are integral and fixed by conjugation, and hence belong to $R$.
Conversely $R+R\theta\subseteq A$. Therefore
\[
 A=R\oplus R\theta,
 \qquad A\simeq R[X]/\bigl(X^2-(\theta+\bar\theta)X
                                  +\theta\bar\theta\bigr).
 \tag{DR.4}
\]
The coefficients lie in $R$ and the discriminant is
$\delta^2\in R^\times$. Reduction modulo a uniformizer $\varpi$
of $R$ is consequently a separable quadratic algebra over the finite
residue field $k$. It is either $k\times k$, with conjugation
exchanging the factors, or the quadratic field $k_2/k$, with its
nontrivial field involution.

In the first case the two simple residue roots lift to roots in $R$.
Explicitly, if $f(a_m)\in\varpi^mR$ and $f'(a_m)$ is a unit, choose
$t_m\in R$ so that
$f(a_m)/\varpi^m+f'(a_m)t_m\equiv0\pmod\varpi$ and put
$a_{m+1}=a_m+\varpi^m t_m$. The errors improve by one power and
the sequence converges to a root. The two lifted roots have unit
difference, so the Chinese remainder theorem in (DR.4) gives
$A\simeq R\times R$. In the second case $A/\varpi A=k_2$ is a
field. An element with nonzero residue is a unit: lift its residue
inverse and invert the resulting $1+\varpi z$ by a convergent
geometric series. Every nonzero element is then a power of
$\varpi$ times a unit, using the finite free $R$-module (DR.4).
Thus $A$ is the unramified quadratic discrete valuation ring over $R$.
This establishes the split/unramified alternatives from the literal
source hypothesis, without assuming them separately.

We need the norm on units in the nonsplit case. Let $|k|=q$.
The finite-field norm $k_2^\times\to k^\times$ is onto, since
it is $z\mapsto z^{q+1}$ on the cyclic group of order $q^2-1$.
The finite-field trace is onto as well. Indeed
$z\mapsto z+z^q$ is a $k$-linear map into the one-dimensional
space $k$, and it is not zero: the nonzero polynomial $X+X^q$
has degree $q<q^2$ and cannot vanish at every element of $k_2$.
This proof works when $\operatorname{char}k=2$.

Let $u\in R^\times$. Choose $b_1\in A^\times$ with
$b_1\bar b_1\equiv u\pmod\varpi$ using the finite-field norm.
Suppose $b_m\bar b_m-u=\varpi^m c_m$. Choose $z_m\in A$
whose residue trace is $-c_m/u$, and set
\[
                    b_{m+1}=b_m(1+\varpi^m z_m).
 \tag{DR.5}
\]
For $m\ge1$ the identity
\[
 (1+\varpi^m z)(1+\varpi^m\bar z)
       \equiv1+\varpi^m(z+\bar z)\pmod{\varpi^{m+1}}
 \tag{DR.6}
\]
holds because $2m\ge m+1$. Hence
$b_{m+1}\bar b_{m+1}\equiv u\pmod{\varpi^{m+1}}$.
The sequence is Cauchy and remains in the same nonzero residue class.
Its limit $b\in A^\times$ satisfies $b\bar b=u$.
No division by $2$ or assumption of odd residue characteristic is
used. In the split case the norm is $(a,b)\mapsto ab$ and is
plainly onto the units of $R$.

Now let $M$ be a free rank-$n$ lattice over the nonsplit ring $A$
with a perfect integral Hermitian pairing, conjugate linear in the
first argument. Its Gram determinant in an integral basis is a unit,
so its reduction is a nondegenerate Hermitian form over $k_2/k$.
If $n>0$, this reduced form has a vector of nonzero self-pairing.
For otherwise, for every $x,y$ and every $a\in k_2$,
\[
 0=h(x+ay,x+ay)=a\,h(x,y)+\bar a\,\overline{h(x,y)}.
 \tag{DR.7}
\]
Nondegeneracy supplies $x,y$ with $h(x,y)\ne0$. As $a$ varies,
the right side is the trace of every element of $k_2$, contradicting
its surjectivity. In particular this argument remains valid in
characteristic two, where the trace of $1$ itself can be zero.

Lift such a vector to $v\in M$. Its norm $h(v,v)=u$ is a unit
of $R$. By (DR.5) choose $b\in A^\times$ with $b\bar b=u$,
and replace $v$ by $b^{-1}v$. It now has norm one. The projection
$x\mapsto v h(v,x)$ gives the integral direct sum
\[
             M=Av\oplus M_1,\qquad M_1=v^\perp\cap M.
 \tag{DR.8}
\]
The complement is free of rank $n-1$, since $A$ is a discrete
valuation ring. It is again perfect: a vector of its fraction-field
span pairing integrally with $M_1$ pairs integrally with all of $M$
and is therefore in $M$, hence in $M_1$. Induction gives an integral
orthonormal basis of $M$. The rank-zero case is immediate.

For the split ring $A=R\times R$, write $M=M_+\oplus M_-$.
The Hermitian pairing is a perfect duality between the two free
$R$-lattices. Choose a basis $e_i$ of $M_+$ and its dual basis
$f_i$ of $M_-$. The vectors $(e_i,f_i)$ form an $A$-basis of $M$
with Gram matrix $I$. Thus the same orthonormal-basis conclusion
holds in both local alternatives. This proves the normal-form
ingredient used by \texttt{unimodGenus}, with the actual completed
CM algebra and its actual involution.

\paragraph{DR2: an actual finite-adelic matrix carrying a self-dual lattice.}
Let $N\subset E^n$ be a supplied full $\mathcal O$-lattice with
$N=N^\#$ for the standard Hermitian pairing. Choose a nonzero
rational integer $d$ with
\[
                 d\mathcal O^n\subseteq N
                         \subseteq d^{-1}\mathcal O^n.
 \tag{DR.9}
\]
Such a $d$ follows by clearing the finitely many coordinates of
generators and of an $E$-basis contained in $N$. Thus $N_P$ is
standard outside the finitely many primes dividing $d$.

At every finite prime $v$ of $F$, let $\widehat N_v$ be the product
of its completed local lattices above $v$. It is a free rank-$n$
module over the integral algebra $A$ in (DR.2). Hermitian duality
commutes with these localizations and completions: in a rational
local basis its defining inequalities are that the coordinates
after the conjugate-transpose Gram matrix are integral. Passing
from the local ring to its completion does not change this dual
matrix formula. At a conjugate pair it pairs the two different
components, exactly as in NB1; at a fixed prime it is the ordinary
unramified Hermitian dual. Consequently $\widehat N_v$ is perfect.

If the prime is nonsplit, DR1 gives an orthonormal integral basis
of $\widehat N_v$. Let $g_v\in\mathrm{GL}_n(B)$ be its column
matrix. Then $g_v^*g_v=I$ and $g_v A^n=\widehat N_v$.
At a split pair $P,\bar P$, one can make the choice even more
explicit. Choose any completed basis matrix $B_P$ of $N_P$ and set
\[
              g_P=B_P,\qquad
              g_{\bar P}=\overline{B_P}^{\,-\mathsf T}.
 \tag{DR.10}
\]
The bar transports entries to the conjugate completion.
The equality $N=N^\#$ says exactly that the second matrix is
a basis of the lattice at $\bar P$. Formula (DR.10) is unitary.
At every place where $N$ is standard, choose the identity matrix.

These choices define a matrix $g$ over the actual finite adelic
algebra $E\otimes_F\mathbb A_{F,f}$. To see this directly, expand
every $g_v$ and its inverse in a fixed $F$-basis $1,\theta$ of $E$.
Each coefficient is its identity-matrix coefficient outside finitely
many $v$, so it is an actual finite adele of $F$. Taking the corresponding
two tensor terms constructs every matrix entry in the stated algebra.
Both the entries and those of the inverse are integral outside the
finite set in (DR.9), where they are the identity.
The inverse is the conjugate transpose
componentwise, so $g$ is unitary. The chosen local bases give
\[
          g\widehat{\mathcal O}^{n}
                    =\prod_P\widehat N_P=\widehat N.
 \tag{DR.11}
\]
Here the last equality is the integral adelic-span identity proved
in AD4 using the finite sandwich and Chinese remaindering. It holds
for adelic vectors, not just rational ones. Thus (DR.11) is the
literal \texttt{Carries} assertion.
For rational vectors, intersection with $E^n$ recovers $N$: membership
in a completed local lattice is the nonnegativity of the valuations
of its rational coefficients in a local basis, hence is local
membership; LP1 then applies. The source's
$\operatorname{globLat}(g)$ is therefore precisely $N$.
This proves the actual \texttt{Formalism.L2.Col.UnimodGenus}
existence needed in \texttt{exists\_isPtOf\_cm}: that predicate is
the adelic $\operatorname{globLat}$ existence assertion just proved.
The similarly named \texttt{UnimodGenusLoc} and
\texttt{PI.UnimodGenus} ask for rational isometries of semilocal
lattices and are not asserted here. The direct adelic construction
bypasses that part of the source route and invokes no rational
approximation by unitary matrices.

\paragraph{DR3: realize every prescribed residue and rational frame.}
Let $p$ be a valid framed decorated lattice of rank $n$, with
underlying self-dual lattice $N$. Its selected type primes $S_T$,
auxiliary primes $S_G$, and rational-frame primes $Q$ are pairwise
disjoint finite sets. Assume the source's exact orientation condition
\[
 P\in S_T\cup S_G\cup Q
       \quad\Longrightarrow\quad
 \bar P\notin S_T\cup S_G\cup Q.                    \tag{DR.12}
\]
In particular each selected prime lies in a split conjugate pair;
neither component is selected twice. The valid decoration supplies
a nonzero residue-linear $f_P:N_P\to\mathcal O/P$ at $P\in S_T$,
a residue-basis map $\mathrm{be}_P:N_P\to(\mathcal O/P)^n$ at
$P\in S_G$, and an exact rational local basis matrix at $P\in Q$.

Choose a rational local basis $B_P$ of $N_P$ for each $P$ in the
union. Such a basis exists because a full lattice over the local
discrete valuation ring is free; its vectors lie in $E^n$ and form
an $E$-basis. Write $k_P=\mathcal O/P$.
The functional $f_P$ descends to a nonzero row
$r_P:k_P^n\to k_P$ in the coordinates $B_P$.
Residue-linearity over $\mathcal O_{(P)}$ follows from its
$\mathcal O$-linearity by clearing a denominator outside $P$.
Choose a basis of $\ker r_P$ and then a vector on which $r_P=1$.
Its column matrix $U_P\in\mathrm{GL}_n(k_P)$ satisfies
\[
                         r_PU_P=e_n^{\mathsf T}.    \tag{DR.13}
\]
Lift its entries through the actual quotient map
$\mathcal O\twoheadrightarrow k_P$ to a rational matrix
$\widetilde U_P$. Its determinant has nonzero residue, so the
matrix and its inverse are integral at $P$. Replacing $B_P$ by
$B_P\widetilde U_P$ preserves its local lattice and makes the last
coordinate of its inverse reduce exactly to $f_P$.

At an auxiliary prime, the supplied map $\mathrm{be}_P$ induces
an isomorphism $N_P/PN_P\simeq k_P^n$.
For clarity its global kernel and surjectivity clauses suffice:
clear a denominator $s\notin P$ to put a vector in $N$;
zero residue then puts $sx$ in $PN$, hence $x$ in $PN_P$.
Surjectivity already has global witnesses. Relative to $B_P$ this
isomorphism is a matrix $M_P\in\mathrm{GL}_n(k_P)$.
Lift $M_P^{-1}$ as above and replace $B_P$ by that lifted
right multiple. Its inverse-coordinate reduction is now the exact
given map $\mathrm{be}_P$, with every coordinate retained.
At a prime of $Q$ use the supplied rational frame matrix itself;
the defining framed-decoration condition says that it is a local
basis of $N_P$. Disjointness makes these prescriptions compatible.

Apply DR2 to obtain an initial unitary matrix carrying the standard
lattice to $N$. Apply AD4 once, with its modification set equal to
$S_T\cup S_G\cup Q$ and the just chosen rational local matrices.
Condition (DR.12) is exactly its no-conjugate-pair premise.
It gives a unitary $g$ carrying the very same $N$, whose component
at each selected $P$ is exactly $B_P$, and whose conjugate component
is its conjugate inverse transpose. Thus its pullback on a global
$x\in N$ has at $P$ the integral coordinates $B_P^{-1}x$.

The residue statements just proved are the source's literal
\texttt{HasRes} statements. To check this last point explicitly,
write any rational integral coordinate as $a/s$ with $s\notin P$.
Choose $r\in\mathcal O$ with
$\rho_P(r)=\rho_P(a)\rho_P(s)^{-1}$.
Then $a/s-r\in P\mathcal O_{(P)}$, so its image in the completion
has multiplicative valuation strictly less than one. This is precisely
$\texttt{HasRes}\,P\,(a/s)\,\rho_P(r)$, using the actual quotient
map. We obtain
\[
\begin{array}{ll}
 P\in S_T:&\operatorname{res}_P((g^{-1}x)_n)=f_P(x),\\
 P\in S_G:&\operatorname{res}_P((g^{-1}x)_j)
                                      =\mathrm{be}_P(x)_j
                                      \quad\text{for every }j,\\
 P\in Q:&g_Pe_j=p.\operatorname{fr}(P)(e_j)
                                      \quad\text{for every }j.
\end{array}                                                     \tag{DR.14}
\]
Together with unitarity and (DR.11), these are all five fields
\texttt{uni}, \texttt{lat}, \texttt{f}, \texttt{be}, and \texttt{fr}
of \texttt{IsPtOf}. If $n=0$, a nonzero functional on the zero
module is impossible, so a valid decoration forces $S_T=\varnothing$.
The empty identity matrix gives the conclusion directly; all auxiliary
and frame coordinate clauses are vacuous.
This proves the full \texttt{exists\_isPtOf\_cm} assertion.

For the true-column application, the source's all-twist wrapper
introduces no further geometry. A type prime contains $T_u$, while
the auxiliary and twist primes do not; these sets and their conjugates
are therefore disjoint. Two selected type primes cannot be conjugate
because $(\pi_u,\bar\pi_u)=\mathcal O$. The supplied twist-level
split condition handles the remaining auxiliary/twist union and its
orientation, while its disjointness clause separates those two sets.
These are exactly the checks in \texttt{exists\_ptOf\_twists}.
Hence DR3 supplies its actual points, the \texttt{yStd} points of FC3,
and the initial realizers used in NB6. AD3 and NB then supply the
supported moves as already proved. No additional point-realizability
hypothesis is needed beyond the literal supplied CM and valid-decoration
data. The full-column terminal incidence and special symmetric-level
identifications, and the upstream existence of the arithmetic host,
remain distinct from this completed realization construction.
